\begin{table}[!htbp]
\centering
\caption{Baseline characteristics of projected-income value-added compliers}
\label{tab:program-income-va-complier-ses}
\begin{threeparttable}
\small
\setlength{\tabcolsep}{4pt}
\begin{tabular}{lcccc}
\toprule
 & SAE mean & Above P25 & Above median & Above P75 \\
\midrule
Income decile & 4.519 & -0.326*** & +0.189*** & +0.749*** \\
 &  & (0.081) & (0.049) & (0.048) \\
Father's education (years) & 10.771 & -0.289*** & +0.165*** & +0.740*** \\
 &  & (0.089) & (0.053) & (0.052) \\
Mother's education (years) & 10.941 & -0.179** & +0.271*** & +0.680*** \\
 &  & (0.084) & (0.050) & (0.049) \\
\bottomrule
\end{tabular}
\begin{tablenotes}[flushleft]
\footnotesize
\item Notes: Each threshold column reports the estimated mean baseline characteristic among students induced by their first-round offer to attend a school above the stated projected-income value-added threshold, minus the corresponding mean among all timely SAE applicants in the 2018--2020 cohorts. The SAE mean is common across thresholds. Parentheses contain heteroskedasticity-robust standard errors that account for estimation of the complier and population means and their covariance. High value added is defined using equally weighted school quantiles of EB-shrunken value added. The complier population and estimation sample may differ across thresholds. Differences describe baseline composition and are not treatment effects. $^{*}p<0.10$, $^{**}p<0.05$, $^{***}p<0.01$.
\end{tablenotes}
\end{threeparttable}
\end{table}
